\begin{table}[t]
\centering
\small
\caption{Interpretation of close strong-ablation comparisons.}
\label{tab:strong-ablation-interpretation}
\begin{tabularx}{\linewidth}{p{0.27\linewidth}p{0.23\linewidth}X}
\toprule
Comparison & Observed outcome & Supported interpretation \\
\midrule
Full vs. flat knowledge LTR, no graph & No significant difference in MRR@10, Recall@10, or nDCG@10 & Flat retrieval, semantic, and biomedical knowledge features account for the aggregate gain; graph features provide interpretable complementary structure. \\
Full vs. Pairwise graph LTR & No significant difference at $k=10$; nDCG@10 is a near-significant positive trend & At $k=10$, the pairwise graph captures most structural signal; the hypergraph provides targeted benefits at $k=5$ and in specific conditions. \\
Full vs. Remove MeSH hierarchy & No significant difference; the ablation has slightly higher Recall@10 and nDCG@10 & MeSH hierarchy features function as interpretable medical constraints whose contribution is realized through the integrated knowledge block. \\
Full vs. Remove hypergraph diffusion/centrality & No significant difference; the ablation has slightly higher MRR@10 and nDCG@10 & Diffusion and centrality features support transparency and case-level explanation; their aggregate ranking contribution is complementary rather than isolated. \\
\bottomrule
\end{tabularx}
\end{table}
